% Limitations and conclusion.

\section{Limitations}
\label{sec:limitations}

The pooling assumption is not testable at this budget in any strong sense. The registered
screen covers \ScreenedModels{} of \NumModels{} models and detects a per-model departure of
\ScreenRho{} with \ScreenPowerPct\% power. Smaller departures, and any departure in the models not
screened, remain outside its reach. This is why the results are reported as a curve over
$\varepsilon$ and why the no-pooling analysis is reported next to them.

The human reference is a careful reader's judgement of what a response commits to, not
mathematical truth. A labeler who transcribes an answer faithfully has not checked whether the
answer is right; correctness is then decided by the benchmark's own equivalence routine
against its own reference answer. Errors in those reference answers propagate to every result
here and to the leaderboard equally.

The three verifiers share a normalisation routine and a symbolic algebra library, so a
response that defeats that shared machinery can be miscounted by all three at once and land in
a concordant cell. Such a response would be invisible to the agreement stratification and
would be caught only by the human labels in the concordant cells, which is one reason the
registered design spends labels there.

The item-sampling term is asymptotic, and the coverage statement says so. A
distribution-free version is available and is substantially wider.

Everything is measured on one benchmark, one leaderboard, and a fixed set of published
generations. The census in Section~\ref{sec:generality} shows how widely the same scoring
design is used; it does not show how often that design changes a ranking elsewhere.

Finally, the audit labels were produced by a small number of labelers, one of whom is the
author. The tool records timing and blocks pasting, batches pass an automated gate before use,
and a subsample is independently re-labelled, but the labels are not crowd-sourced and their
reliability is measured rather than assumed.

\section{Conclusion}
\label{sec:conclusion}

A leaderboard reports one number per model and names only half of what produced it. The other
half is a verifier, and on the board studied here the choice of verifier moves the ranking,
which a published score gives a reader no way to detect.

What a small human audit can do about that is limited and worth stating precisely. It can
certify part of the ranking, but only under an assumption about how models relate to one
another where the verifiers agree; without that assumption, this audit certifies
\CertNoPooling{} of \NumOrderings{} orderings. With it, at the tolerance registered before
seeing any result, it certifies \CertPrereg{}. The useful output is neither number on its own
but the curve between them, together with the breakdown value of each ordering.

The practical recommendation is narrow and cheap. A benchmark that reports how often its
scorer failed to read a response, alongside the score, would let a reader compute the bound in
Section~\ref{sec:setup} without any audit at all. The count already exists inside the scoring
pipeline; it is discarded.
